% !TEX root = ../main.tex

\chapter{面向股票预测的多源信息融合框架设计}

\section{任务定义}

本文研究的是单只股票在给定分析日期下的短期走势判断。给定目标股票、分析日期，以及该日期之前能够看到的行情和文本材料，离线实验需要判断该股票在未来 1 个、3 个和 7 个交易日后的价格方向，即相对分析日收盘价更可能上涨还是下跌。本文没有直接做精确涨跌幅回归，主要是因为短期涨跌幅噪声更大，对当前数据质量和样本规模要求也更高；先做方向判断更适合作为本系统的第一步。

输入材料可以分成两类。一类是 K 线数据，包括开盘价、最高价、最低价、收盘价、成交量、涨跌幅、换手率等字段；另一类是评论、新闻、研报和搜索结果等文本。前者提供价格和量能状态，后者补充事件、情绪和分析观点。本文希望用大语言模型把这两类材料放到同一个判断问题中，而不是只在行情表格上训练一个孤立分类器。

\section{总体框架}

MarketMind 的整体流程可以理解为三步：先围绕目标股票收集材料，再把过长或噪声较高的文本压缩成摘要，最后把摘要和 K 线一起交给模型判断。这里的取舍是，摘要会损失一部分原文细节，但能避免模型直接面对大量重复评论、长篇研报和搜索网页，从而让最终结论更容易追溯到具体材料。

从实现上看，系统仍然可以拆成数据层、信息处理层、模型判断层和输出层。数据层维护行情、评论、新闻和研报等原始数据；信息处理层负责筛选、去重和压缩；模型判断层比较不同来源之间的一致和冲突；输出层则根据使用场景分成两种形式：在线页面输出趋势判断和理由，离线实验保留多周期方向标签以便量化评估。

采用这种分层处理，主要是因为不同材料的长度和噪声差异很大。若把原始评论、新闻和研报一次性拼接给模型，不仅会浪费上下文，也容易让重复信息或低质量文本干扰判断；若压缩过度，又可能丢掉少数关键事件。本文折中采用分支摘要作为中间表示，先保留每类材料的主要观点和风险，再进入最终预测环节。

在当前在线系统中，这一框架已经对应到一条可执行链路：系统根据股票代码生成实时信息快照，其中包括最近 50 条 comments、最近 10 条 news、最近 3 条 reports、最近 7 个交易日 K 线以及 SerpAPI 补充检索结果；用户可以按需对 comments、news、reports 和 SerpAPI 结果触发 Summary；最后，统一预测 Prompt 将 K 线与已生成的 news / reports / SerpAPI Summary 输入模型，输出未来一段时间的整体趋势判断。这样，框架不仅停留在离线方法描述上，也在在线原型系统中形成了“抓取—摘要—综合预测”的闭环。

\section{多源信息来源设计}

\subsection{行情信息}

行情信息是股票方向预测中最基础的数据来源。本文使用目标日期之前连续若干个交易日的 K 线作为结构化输入。对于监督微调主线，系统使用连续 7 个交易日的归一化 K 线预测第 8 个交易日方向；对于多源信息验证，系统使用目标日期及之前连续 5 个交易日的归一化 K 线，并预测未来 1、3、7 个交易日方向。

行情数据的作用主要有两个。第一，它提供价格和成交量层面的短期状态，例如趋势是否延续、是否出现放量反弹、是否存在连续下跌或高位震荡。第二，它为文本信息提供市场背景。相同的新闻或评论在不同技术形态下可能产生不同影响，因此模型需要同时看到行情状态和外部文本，才能形成更合理的短期判断。

\subsection{评论信息}

评论信息主要用于捕捉投资者情绪、讨论热度和短期市场叙事。论坛或股吧中的评论通常更新较快，能够较早反映散户投资者对热点题材、突发消息和价格波动的反应，这与社交媒体和投资社区情绪可作为市场先行信号的相关研究结论基本一致\cite{bollen2011twitter,zhang2025gnn,zhou2025expertopinion}。但评论文本也包含大量情绪宣泄、重复观点和非理性表达。本文保留评论分支，是为了利用它的时效性；不直接把原文作为最终预测依据，则是为了避免情绪噪声压过新闻和研报中的事实信息。

在框架设计中，评论分支的目标不是保留所有原始评论，而是提取其中的代表性观点和情绪共识。系统优先选择互动量较高或关注度较强的评论标题，再引导大语言模型归纳主要观点簇、情绪倾向和共识程度。这样做相当于在“保留市场情绪”和“降低噪声干扰”之间取一个中间方案。

\subsection{新闻信息}

新闻信息主要体现事件驱动特征。公司经营变化、行业政策、市场传闻、公告解读和突发风险通常会通过新闻报道进入市场信息流，并在短期内影响投资者预期。已有研究表明，大语言模型能够从新闻标题和金融新闻文本中提取对未来收益方向具有解释力的关键信息\cite{lopezlira2023chatgptstock,elahi2024financialnewsllm}。与评论相比，新闻文本结构更规范、事实性更强，但也可能存在报道角度差异、标题信息不足和发布时间滞后等问题。

在本文框架中，新闻分支主要用于提取与未来走势相关的事件线索。系统将目标日期之前可见的相关新闻组织为标题或摘要列表，要求模型归纳主要事件主题、媒体叙述倾向以及可能影响股价的关键信息。对于更完整的后续系统，新闻分支还可以采用两阶段方式：先基于标题筛选重要新闻，再读取正文进行二次摘要。

\subsection{研报信息}

研报信息主要提供专业分析和中长期逻辑。与评论和新闻相比，研报通常包含更完整的行业背景、公司基本面分析、业绩预期、风险提示和投资逻辑。金融大语言模型相关综述普遍指出，金融文本理解不仅依赖通用语言能力，也依赖对领域术语、分析逻辑和上下文约束的把握\cite{lee2024finllms,li2023llmfinance,nie2024llmfinancialapplications}。虽然研报的时效性不一定最强，但其分析深度能够帮助模型补充更稳定的解释框架。

在框架设计中，研报分支用于提供高质量分析证据。系统选取分析日期之前较近的研报，对其核心观点、逻辑主线、潜在机会和风险因素进行摘要。最终预测时，研报摘要可以与新闻事件和评论情绪共同构成跨源证据，使模型不只依赖短期情绪，也能够参考相对专业的分析观点。

\subsection{Web API 检索信息}

除 comments、news 和 reports 外，本文进一步将外部 Web API 检索结果作为补充信息源接入框架。该部分的目标不是替代已有金融站点抓取，而是在实时查询阶段为目标股票补充更广泛的公开网页信息，例如最新媒体报道、公告解读、股价异动讨论、财报相关新闻以及面向个股的研究观点等。这种“在生成前先接入外部可更新知识”的思路与检索增强生成的基本方向一致\cite{lewis2020rag,ye2024r2ag}。

在当前实现中，Web API 检索通过 SerpAPI 调用搜索引擎结果完成。系统围绕目标股票构造多组检索 query，包括“公司名 + 股票 + 最新消息”“公司名 + 股价 + 异动”“公司名 + 公告”“公司名 + 财报/业绩”“公司名 + 研报”以及“股票代码 + 公司名”等形式。每组 query 默认获取前 3 条结果，并对返回内容中的标题、链接、摘要、来源和日期进行抽取。为了减少不同检索意图之间的大量重复网页，系统在多组 query 返回后进一步以 URL 为主键进行统一去重与合并，最终形成一份补充搜索结果列表。

与新闻分支相比，Web API 检索的覆盖范围更广，不局限于单一资讯站点，适合补充临时热点、媒体扩散信息和交叉验证证据。但它的噪声也更高，不同 query 之间容易出现重复结果、弱相关结果或摘要质量不稳定的问题。因此，Web API 检索结果在本文框架中不替代 news/report 分支，而是作为最后一层补充证据，帮助模型判断是否存在新的外部事件线索或证据支持。检索结果进入模型前仍需要较严格的组织与筛选，以减少“检索到信息”和“模型真正利用信息”之间的落差\cite{ye2024r2ag}。

\section{多源信息组织与 Prompt 设计}

为了让模型稳定理解输入，本文在 Prompt 中固定展示股票代码、分析日期、行情窗口、信息来源和对应摘要。不同来源分块呈现，避免评论、新闻和研报混在一起，使模型难以判断某条信息来自哪里。

各分支 Prompt 的侧重点并不相同。评论分支更关注讨论主题、整体多空倾向和争议点；新闻分支更关注近期发生了什么，以及这些事件可能带来的正负影响；研报分支偏向分析逻辑、业绩预期和风险提示；Web API 检索分支则保留来源、日期和摘要，帮助模型判断搜索结果是否提供了新的线索。综合预测 Prompt 在此基础上比较不同来源之间的一致与冲突，并结合 K 线状态输出未来趋势判断。

在当前实现中，Prompt 分为两类。第一类是分支 Summary Prompt，分别针对 comments、news、reports 和 SerpAPI 网页正文生成集中总结。comments 的输入以标题、发布时间、点击量和评论量为主，用来观察高互动讨论中的多空倾向和争议；news 与 reports 的输入保留标题、日期和正文内容，用来提取事件线索、分析逻辑和风险提示；SerpAPI 部分会先经过规则过滤，剔除 PDF 链接，再对剩余普通网页抓取正文并总结，同时提示模型注意来源可靠性、重复信息和弱相关内容。第二类是综合预测 Prompt，它不再直接消费所有原始文本，而是只接收最近 7 个交易日 K 线以及已生成的 news / reports / SerpAPI Summary，并据此形成未来趋势判断。这样的设计把 comments 保留为辅助观察分支，把更偏事实和分析性的三个摘要分支作为最终预测的主要外部材料。

当前在线原型系统中的综合预测 Prompt 不再要求模型固定输出未来 1、3、7 个交易日的涨跌方向，而是要求模型在阅读 K 线和各类 Summary 后，给出更接近投研分析的未来趋势预测。输出结构包括“趋势预测”“主要利好/利多依据”“主要利空或风险”“关键不确定性”和“综合结论”五部分。相比三项机械标签，这种形式更适合在线演示，也能保留判断依据。

这种 Prompt 设计的关键在于“先压缩、再融合”。分支摘要阶段减少原始文本长度和噪声，综合预测阶段再把不同来源的信息放到同一个决策问题中。相比简单拼接原始文本，该方法更容易控制上下文长度，也更容易让模型给出可解释的预测依据。这与检索增强和 retrieval-aware 组织方式中“先组织外部证据、再引导模型利用证据”的思路一致\cite{lewis2020rag,ye2024r2ag}。

\section{预测输出设计}

本文将预测输出区分为离线评测口径和在线展示口径。离线实验为了便于量化比较，仍采用未来 1 日、3 日和 7 日的方向标签；在线原型系统更强调辅助分析价值，不再强制输出固定多周期涨跌标签，而是输出对股票未来趋势的综合判断，并说明判断依据主要来自哪些 K 线特征、新闻事件、研报逻辑或外部搜索证据。

离线实验中，模型输出会被解析为“涨”或“跌”标签，用于计算准确率和分周期表现。在线应用更强调解释依据：用户不仅关心预测结果，也需要判断该结果是否来自可靠证据、是否存在信息源冲突、是否还需要人工复核。当前网站的 Predict 按钮按这一思路实现，先由用户按需生成各分支 Summary，再在统一预测入口输出趋势预测、上涨依据、下行风险和不确定因素，把原始信息浏览、分支总结和综合判断组织为一套交互流程。

\section{本章小结}

本章把 MarketMind 的预测流程拆成了几个比较具体的环节：先抓取行情、评论、新闻、研报和搜索结果，再分别做筛选与摘要，最后把 K 线和摘要送入综合预测 Prompt。这样做的主要考虑是控制上下文长度和噪声，而不是简单把所有原文拼接给模型。在线系统最后输出的是趋势判断和依据；离线实验为了计算指标，仍保留未来 1、3、7 日方向标签。下一章会继续说明这些输入如何被整理成训练样本和评测样本。
